\chapter{Calibration tables}
\label{app:calibration-tables}

\begin{table}[H]
  \centering
  \begin{tabular}{lr}
    \toprule
    \textbf{Filtering step} & \textbf{Observations} \\
    \midrule
    Initial observations & 3,778 \\
    Observations after cleaning & 3,665 \\
    \midrule
    Removed ITM options & 1,889 \\
    \quad Original call observations & 417 \\
    \quad Puts converted to calls & 1,359 \\
    Total call observations & 1,776 \\
    \midrule
    Removed by price bounds & 0 \\
    Removed by slope/convexity & 486 \\
    \quad Monotonicity violations found & 12 \\
    \quad Convexity violations found & 7,449 \\
    Removed by maturity filter & 0 \\
    \midrule
    \textbf{Remaining observations} & \textbf{1,290} \\
    Remaining maturities & 12 \\
    \bottomrule
  \end{tabular}
  \caption{Summary of option data filtering.}
  \label{tab:option-filtering}
\end{table}

\begin{table}[H]
    \centering
    \begin{tabular}{lr}
        \toprule
        Maturity & Discount Rate (\%) \\
        \midrule
        2019-12-20 & 1.6367 \\
        2020-01-17 & 2.2623 \\
        2020-02-21 & 2.1119 \\
        2020-03-20 & 2.0381 \\
        2020-04-17 & 2.0846 \\
        2020-06-19 & 1.9729 \\
        2020-09-18 & 1.9206 \\
        2020-10-16 & 1.9557 \\
        2020-11-20 & 1.9104 \\
        2020-12-18 & 1.8774 \\
        2021-06-18 & 1.8524 \\
        2021-12-17 & 1.8363 \\
        \midrule
        Inferred dividend yield & 1.8914 \\
        \bottomrule
    \end{tabular}
    \caption{Inferred discount rates by maturity}
    \label{tab:inferred_r_q}
\end{table}

\begin{table}[H]
  \centering
  \begin{tabular}{lcccccc}
    \toprule
    Weight Scheme & Model & IV RMSE &IV MAE &IV Bias& W.RMS & Parameters at bounds\\
    \midrule
    SSE & SVJJ & 1.122 &0.661 & 0.186& 2.048 & $\sigma$ \\
    SSE & SVJ & 1.268 & 0.691& 0.272& 2.068 & None \\
    SSE & Heston & 1.203 &0.784&-0.091& 2.351 & None \\
    \midrule
    Bid-Ask& SVJJ & 1.213 & 0.588&-0.270&2.689 & $\rho, \overline{v}$ \\
    Bid-Ask& SVJ & 0.994 &0.504&-0.021& 2.859 & None \\
    Bid-Ask& Heston & 1.020 &0.601&-0.108& 3.376 & None \\
    \midrule
    Vega & SVJJ & 0.953 &0.452&-0.148& 0.439 & None \\
    Vega & SVJ & 0.955 &0.467&-0.142& 0.455 & None \\
    Vega & Heston & 0.993 &0.578&-0.211& 0.585& None \\
    \bottomrule
  \end{tabular}
  \caption{Calibration performance for the different weighting schemes.
    IV figures are measured in volatility points.
    W.RMS denotes the weighted root mean square residual, and is measured in
    price units for SSE and Bid-Ask and in volatility points
  for Vega.}
  \label{tab:calibration-weightsheme}
\end{table}


\begin{table}[ht]
\centering
\begin{tabular}{lrlr}
\toprule
Parameter & Value & Parameter & Value \\
\midrule
$\lambda$  & $1.2098$  & $\sigma$ & $0.0137$ \\
$\rho$     & $-0.8435$ & $\rho_j$ & $-0.7880$ \\
$v$        & $0.0125$  & $\mu_v$  & $0.0651$ \\
$\kappa_v$ & $2.7564$  & $V_0$    & $0.0117$ \\
$\sigma_v$ & $0.4748$  & $\mu_s$  & $0.0227$ \\
\bottomrule
\end{tabular}
\caption{Calibrated SVJJ parameters (Vega-weighted objective).}
\label{tab:svjj-params}
\end{table}


\global\def\lastproofcolor{formulabordercolor}
\chapter{Proofs}
\section{Proof of \texorpdfstring{\Cref{prop:svjj-mg-cond}}{SVJJ MG Cond}}
\label{app:proof-mg-cond}
  By the Itô product rule, the dynamics of $M_t$ are:
  \[
    \frac{\d M_t}{M_{t-}} = -\overline{\mu}\lambda\d t +\sqrt{V_t}\d B_t +(\ee^{J^S}-1)\d N_t.
  \]
  Since the jumps are driven by a marked Poisson process, their expected contribution to the drift
  per unit of time is, the expected contribution from the jumps per unit of time, i.e.
  $\lambda\E[\ee^{J^S}-1]\d t$. The drift of $M_t$ is therefore $\lambda(
  \mathbb{E}[\ee^{J^S}-1]-\overline{\mu})\d t.$ Under $\Q$ the discounted stock price is a
  martingale, hence the martingale condition implies
  \[ 
    \overline{\mu}=\mathbb{E}[\ee^{J^S}-1].
  \]
  Using the conditional distribution of $J^S$, we obtain
  \[
    \mathbb{E}[\ee^{J^S}]
    =\mathbb{E}\left[ \mathbb{E}[\ee^{J^S}\mid J^V] \right]
    =\ee^{\mu_s+\frac{1}{2}\sigma^2} \mathbb{E}[\ee^{\rho_jJ^V}]
    =\frac{\ee^{\mu_s+\frac{1}{2}\sigma^2}} {1-\rho_j\mu_v},
  \]
  where the last equality uses the moment generating function of $J^{V} \sim \exp(\mu_v)$.
  Subtracting one gives the stated result.
\qed
\section{Proof of \texorpdfstring{\Cref{lem:svjj-jump-trans}}{SVJJ Jump transform}}
\label{app:proof-jump-trans}
By definition, the jump transform is
\[
\phi(c)=\E\left[\exp(c_1J^S+c_2J^V)\right].
\]
Conditioning on \(J^V\) and using the conditional normal distribution of \(J^S\), we obtain
\begin{align*}
\phi(c)
&=\E\left[ e^{c_2J^V} \E\left[\ee^{c_1J^S}\mid J^V\right] \right]
&=\ee^{\mu_sc_1+\frac12\sigma^2c_1^2} \E\left[\ee^{(c_2+\rho_jc_1)J^V}\right].
\end{align*}
Since $J^V\sim\exp(\mu_v)$, its moment-generating function gives
\[
\E\left[\ee^{(c_2+\rho_jc_1)J^V}\right] = \frac{1}{1-\mu_v(c_2+\rho_jc_1)}.
\]
Combining these expressions yields \eqref{eq:svjj-jump-trans}. 
\qed

